% SOURCE: written LAST, from the contribution stack pinned in main.tex and
% related_work.md §1, §2, §8. Every number already appears in the body.
% Citations are anchors already used in Three Tiers and Related Work; none
% is new and refs.bib is unchanged at 31 keys.
% FIREWALL: honest frame stated UP FRONT, not deferred to Limitations. The
% three occupied territories are disclaimed in the same passage as the three
% claims. No banned phrase, no clinical-validity claim, generic risk only.
\section{Introduction}
\label{sec:intro}

A risk score computed at a rural antenatal visit is read by at least three
people. A clinician wants the decomposition and can argue with it. A
community health worker wants to know whether to refer this mother and what
to write down. The mother wants to know what to do next, and she may not read
at all, in which case the score reaches her as a colour and a spoken sentence
or it does not reach her.

Most work on explaining a model to several audiences varies how much of the
explanation each one gets. That is the right instinct, and it leaves one
question open. When the model is not confident, what is each reader
\emph{entitled to be told}?

Consider a case from the dataset we use. The model puts it at $0.013$,
$0.487$ and $0.499$ across low, mid and high risk. The top two are separated
by $0.012$, and the model is wrong: the record is true mid-risk. A clinician
shown that split can see the call was nearly a tie. The mother, at the end of
the same pipeline, has no way to discount anything, and if the rendering
follows the argmax alone she is the one who receives an emergency generated
by a coin flip.

Our response is to make the abstention decision part of the rendering rather
than part of the model. One band is computed per prediction from the class
probabilities, and the highest-severity signal is raised only when the model
is not on a knife edge. Below the threshold the mother's channel is derated
while the clinician's is untouched. \textbf{The derate is red to amber, never
red to green.} Amber still says \emph{needs follow-up} and still sends her to
a health worker. Nothing in this design removes a mother from care, and a
reader who takes ``not an emergency'' to mean ``not a concern'' has read it
backwards.

\subsection{What we claim, and what we do not}

We claim three things.

First, and this is the headline, \textbf{abstention differentiated by
stakeholder vulnerability}. The decision to withhold is taken once, from the
model's output, then rendered as different assertions for different readers,
graded so that the reader with the least power to interrogate the system
receives the most conservative signal.

Second, \textbf{deterministic templating over SHAP attributions}. Every
string a health worker or a mother receives is a fixed template filled from
the case's attributions. No language model runs at inference, so the
rendering can be regenerated, diffed and audited, and there is nothing in it
to hallucinate.

Third, \textbf{a quantified leakage audit} of a widely used benchmark. The
UCI Maternal Health Risk dataset contains 561 exact duplicate rows. Scoring
the same model under the same cross-validation with those rows retained moves
macro-recall from $0.580$ to $0.859$, which lands inside the accuracy band
the literature reports for this dataset.

The disclaimers belong in the same breath. \textbf{We do not claim that
explanations should be audience-specific}, which is settled and was settled
before us~\citep{arya2019explanation,imrie2023multiple,kim2024stakeholder}.
\textbf{We do not claim the gating mechanism.} Declining to answer under
uncertainty has a long lineage in selective prediction and learning to defer,
and our fixed-margin rule is a hand-set and cruder instance of
it~\citep{hendrickx2024reject,mozannar2020consistent}. \textbf{We do not
claim a modelling contribution.} Fitting a classifier with SHAP attributions
to this dataset is saturated~\citep{mamun2025identification,
maisoon2026stratification}, and our model is a vehicle for the rendering
question. On the audit we claim the measurement and not the discovery: the
duplicates are the kind of thing noticed in a notebook and never written up,
and we report what they cost rather than asserting we found them first.

\subsection{What kind of evidence this is}

This is a design and feasibility artifact, and we would rather say so here
than let a reader discover it in Section~\ref{sec:limitations}. No mother, no
community health worker and no clinician has used any of these renderings.

Our evaluation is a formative, single-evaluator heuristic critique conducted
on our own artifact. It found five defects, including a driver list that
could print a reading arguing against its own prediction under a heading
announcing it as a flag, and all five were fixed before this paper was
written. That is what the evaluation is: an instrument that improved the
design. It is not evidence that any tier helps anyone, and the nearest
empirical work argues against assuming so, since community health workers in
India have been reported struggling with AI explanations even after
simplification~\citep{okolo2024easy,solanokamaiko2024explorable}. Studies
with the three groups are the next step, gated on ethics clearance. The data
here is public, the task is generic maternal-health risk, and nothing in this
paper predicts an outcome for a real patient.

Section~\ref{sec:related} places the work against the two literatures it sits
between, Section~\ref{sec:data} covers the dataset and the leakage
measurement, and Section~\ref{sec:tiers} gives the rendering rule and the
three tiers. Section~\ref{sec:eval} reports the evaluation defects first,
Section~\ref{sec:limitations} collects what the design cannot yet support,
and Section~\ref{sec:repro} gives the commands that regenerate every artifact
here.
